\documentclass{stex}

\libusepackage{naproche}
\libusepackage{copyright}
\libinput{preamble}

\title{The Knaster-Tarski Fixed Point Theorem in \Naproche}
\author{Alexander Holz, Marcel Schütz}
\date{2025}

\begin{document}
\begin{smodule}{knaster-tarski.ftl}

\maketitle

\importmodule[libraries/foundations]{maps-and-systems-of-sets.ftl}
\importmodule[libraries/set-theory]{zfc.ftl}

\symdecl*{Knaster-Tarski}

\noindent This is a formalization of the \textit{Knaster-Tarski Fixed Point
Theorem}, i.e. of the assertion that every subset-preserving map has a fixed
point \cite{Schroeder2012}.

\begin{theorem}[forthel,title=Knaster-Tarski,name=Knaster-Tarski]
  Let $x$ be a set.
  Let $h$ be a map from $\Pow(x)$ to $\Pow(x)$ that preserves subsets.
  Then $h$ has a fixed point.
\end{theorem}
\begin{proof}[forthel]
  (1) Define $A \DefEq \SClass{y}{\Pow(x)}{y \Subclass h(y)}$.
  Then $A$ is a subset of $\Pow(x)$.
  We have $\UnionOver A \Elem \Pow(x)$.

  Let us show that (2) $\UnionOver A \Subclass h(\UnionOver A)$.
    Let $u \Elem \UnionOver A$.
    Take $y \Elem A$ such that $u \Elem y$.
    Then $u \Elem h(y)$.
    We have $y \Subclass \UnionOver A$.
    Hence $h(y) \Subclass h(\UnionOver A)$.
    Indeed $h$ is a subset preserving map between systems of sets and $y, \UnionOver A \Elem \Dom(h)$.
    Thus $h(y) \Subclass h(\UnionOver A)$.
    Therefore $u \Elem h(\UnionOver A)$.
  End.

  Then $h(\UnionOver A) \Elem A$ (by 1).
  Indeed $h(\UnionOver A) \Subclass x$.
  (3) Hence $h(\UnionOver A) \Subclass \UnionOver A$.
  Indeed every element of $h(\UnionOver A)$ is an element of some element of
  $A$.

  Thus $h(\UnionOver A) \Eq \UnionOver A$ (by 2, 3).
\end{proof}


\printbibliography
\printlicense[CcByNcSa]{Alexander Holz (2018), Marcel Schütz (2018--2025)}

\end{smodule}
\end{document}
